% ============================================================
%  Abstract (ภาษาอังกฤษ + ภาษาไทย)
% ============================================================

% ── English Abstract ─────────────────────────────────────────
\clearpage
\phantomsection
\addcontentsline{toc}{chapter}{Abstract}

\begin{tabular}{@{}lp{10.5cm}@{}}
  Project Title   & An Exploration of Artificial Intelligence for Detection of \\
                  & Inhibition Zones in Antibiotic Susceptibility Testing \\[0.3cm]
  Credits         & 3 \\[0.2cm]
  Member(s)       & Mr. Peeraphat Siriphatthanakhun \\
                  & Mr. Phattaradon Thanaprasitpattana \\[0.3cm]
  Project Advisor & [TBD] \\[0.2cm]
  Co-advisor      & [TBD] \\
                  & [TBD] \\
                  & [TBD] \\
                  & [TBD] \\[0.3cm]
  Program         & Bachelor of Engineering \\[0.1cm]
  Field of Study  & Health Data Science \\[0.1cm]
  Department      & Health Data Science \\[0.1cm]
  Faculty         & Engineering \\[0.1cm]
  Academic Year   & 2025 \\
\end{tabular}

\vspace{1cm}
\begin{center}\textbf{Abstract}\end{center}

\hspace{\parindent}%
Antibiotic resistance is a critical global health challenge that demands accurate and rapid diagnostic methods.
The Kirby-Bauer disk diffusion method is the internationally accepted standard for Antimicrobial Susceptibility Testing (AST),
yet manual measurement of zone of inhibition (ZOI) diameters remains time-consuming and prone to human error.
This project develops an artificial intelligence (AI)-based web application to automate the detection
and measurement of inhibition zones from photographs of antibiotic susceptibility testing plates.

The system integrates deep learning object detection models---Faster R-CNN with a ResNet-50 backbone and
YOLOv11n---together with classical image processing via Hough Circle Transform, to localize and measure ZOI diameters.
An Optical Character Recognition (OCR) pipeline using EasyOCR with preprocessing (global thresholding,
adaptive binarization, and image rotation) identifies antibiotic disk labels on the plate.
Detected zone diameters are interpreted against CLSI 2025 and EUCAST 2023 breakpoints
to classify bacteria as Susceptible (S), Intermediate (I), or Resistant (R).

A dataset of 287 real laboratory images was augmented to 543 training images and supplemented with 2,000 synthetic images
to improve model generalization.
Experimental results demonstrate that both deep learning models outperform Hough Transform in ZOI detection accuracy.
YOLOv11n achieved superior accuracy with a mean absolute error (MAE) of approximately \textbf{2.5--2.7\,mm}
and a mAP@0.5 of 95.91\%, significantly outperforming the Faster R-CNN baseline (MAE 4.65\,mm).
The complete system, including the React frontend, FastAPI backend, and SQLite database, is fully integrated
and has been \textbf{successfully deployed for production use at the Princess Srisavangavadhana School of Medicine, Chulabhorn Royal Academy} hospital server.
User evaluation by five domain experts yielded an overall usefulness score of 3.60 out of 5.00,
confirming the system's readiness for real-world microbiological laboratory workflows.

\noindent\textbf{Keywords:} Antimicrobial Susceptibility Testing, Zone of Inhibition, Deep Learning, Object Detection,
Faster R-CNN, YOLOv11, Hough Transform, EasyOCR, Image Processing, CLSI, EUCAST, Web Application

\newpage

% ── บทคัดย่อ ─────────────────────────────────────────────────
\clearpage
\phantomsection
\addcontentsline{toc}{chapter}{บทคัดย่อ}

\begin{tabular}{@{}lp{10cm}@{}}
  หัวข้อปริญญานิพนธ์  & การสำรวจปัญญาประดิษฐ์สำหรับการตรวจจับโซนการยับยั้ง \\
                      & ในการทดสอบความไวต่อยาปฏิชีวนะ \\[0.3cm]
  หน่วยกิต            & 3 \\[0.2cm]
  ผู้เขียน             & นายพีรพัฒน์ ศิริพัฒนคุณ \\
                      & นายภัทรดนย์ ธนประสิทธิ์พัฒนา \\[0.3cm]
  อาจารย์ที่ปรึกษา     & ผศ.\ ดร.\ นนทพรรษน์ ลีราช \\[0.2cm]
  ที่ปรึกษาร่วม        & ผศ.\ ดร.\ ปณิธิ อัจฉราฤทธิ์ \\
                      & ดร.\ คริษฐา แจ้งสามสี \\
                      & ดร.\ จิรพันธ์ เปรมสุริยา \\
                      & ดร.\ จุมพล พลวิชัย \\[0.3cm]
  หลักสูตร             & วิทยาศาสตรบัณฑิต \\[0.1cm]
  สาขาวิชา            & วิทยาศาสตร์ข้อมูลสุขภาพ \\[0.1cm]
  ภาควิชา             & วิทยาศาสตร์ข้อมูลสุขภาพ \\[0.1cm]
  คณะ                 & วิศวกรรมศาสตร์ \\[0.1cm]
  ปีการศึกษา          & 2568 \\
\end{tabular}

\vspace{1cm}
\begin{center}\textbf{บทคัดย่อ}\end{center}

\hspace{\parindent}%
ปัญหาการดื้อยาต้านจุลชีพ (Antibiotic Resistance) เป็นความท้าทายระดับโลกที่ต้องการวิธีการวินิจฉัยที่รวดเร็วและแม่นยำ
การทดสอบความไวของเชื้อแบคทีเรียต่อยาปฏิชีวนะด้วยวิธี Kirby-Bauer disk diffusion เป็นมาตรฐานสากลที่ได้รับการยอมรับในระดับนานาชาติ
อย่างไรก็ตาม การวัดขนาดเส้นผ่านศูนย์กลางของโซนการยับยั้ง (Zone of Inhibition: ZOI) ด้วยตนเองนั้นยังคงใช้เวลานานและเสี่ยงต่อความผิดพลาดจากตัวบุคคล (Human Error)
โครงงานนี้จึงมุ่งพัฒนาเว็บแอปพลิเคชันที่ใช้ปัญญาประดิษฐ์ (AI) เพื่อตรวจจับและวัดขนาดโซนการยับยั้งจากภาพถ่ายจานเพาะเชื้อโดยอัตโนมัติ

ระบบที่พัฒนาขึ้นได้รวมโมเดลตรวจจับวัตถุแบบ Deep Learning ได้แก่ YOLOv11n (โมเดลหลัก) และ Faster R-CNN
ร่วมกับการประมวลผลภาพดิจิทัล (Image Processing) เพื่อระบุตำแหน่งและวัดขนาด ZOI
นอกจากนี้ยังใช้ระบบ Optical Character Recognition (OCR) โดย EasyOCR ร่วมกับเทคนิค Preprocessing
(Global Thresholding, Adaptive Binarization และ Image Rotation) เพื่อระบุชื่อยาปฏิชีวนะจากรหัสบนแผ่นยา
ผลการวัดจะถูกนำไปแปลผลตามมาตรฐาน CLSI 2025 และ EUCAST 2023 เพื่อจำแนกเชื้อแบคทีเรียเป็นกลุ่ม
ไวต่อยา (Susceptible: S), ระดับปานกลาง (Intermediate: I) หรือ ดื้อยา (Resistant: R) โดยอัตโนมัติ

ในการทดลองได้ใช้ชุดข้อมูลภาพจริงจากห้องปฏิบัติการ 287 ภาพ และเพิ่มจำนวนเป็น 543 ภาพผ่านเทคนิค Augmentation
ร่วมกับภาพจำลอง (Synthetic Images) อีก 2,000 ภาพ เพื่อเพิ่มความสามารถในการเรียนรู้ของโมเดล
ผลการทดสอบพบว่า YOLOv11n มีประสิทธิภาพสูงสุด โดยมีค่าความคลาดเคลื่อนเฉลี่ย (Mean Absolute Error: MAE)
อยู่ที่ประมาณ \textbf{2.5--2.7 มม.} และมีค่า mAP@0.5 สูงถึง 95.91\%
ซึ่งดีกว่าโมเดล Faster R-CNN และวิธี Hough Transform อย่างมีนัยสำคัญ
ปัจจุบันระบบได้รับการพัฒนาเป็นเว็บแอปพลิเคชันแบบ Full-stack และติดตั้งใช้งานจริง (Production Deployment)
บนเครื่องเซิร์ฟเวอร์ของ \textbf{ราชวิทยาลัยจุฬาภรณ์ วิทยาลัยแพทยศาสตร์ศรีสวางควัฒน}
โดยผลการประเมินจากผู้เชี่ยวชาญ 5 ท่าน พบว่าระบบมีความเหมาะสมในการใช้งานจริงในระดับ 3.60 จาก 5.00 คะแนน
ซึ่งยืนยันถึงความพร้อมในการสนับสนุนการทำงานในห้องปฏิบัติการจุลชีววิทยาคลินิก

\noindent\textbf{คำสำคัญ:} การทดสอบความไวต่อยาต้านจุลชีพ, โซนการยับยั้ง, Deep Learning, การตรวจจับวัตถุ,
YOLOv11, การประมวลผลภาพ, มาตรฐาน CLSI, มาตรฐาน EUCAST, เว็บแอปพลิเคชัน

\newpage
